\section{Discussion}
\subsection{Reframing}
\label{sec:reframing}
The pre-registered H1 criterion was not demonstrated. H1 required a per-benchmark Spearman of at least 0.6 and a win over both StatNN and the norm-of-delta baseline. No learned row reaches 0.6 on either target. Across the 25 complete rows, the highest three-seed mean on HellaSwag is 0.46, from M1-matched on the trajectory split. The highest on MMLU-Pro is 0.45, shared by M1-matched on the trajectory split and M2-matched on the seed split. H1 also has a falsification rule: Spearman below 0.3 on all benchmarks after honest tuning. Several rows exceed 0.3, so that rule was not triggered. H1 is unresolved rather than refuted.

On the seed and trajectory splits, the three-seed mean of every learned row is below seed-0 StatNN on both targets. On the family split, StatNN is near zero, and two learned rows, M1-matched and M6, exceed it on MMLU-Pro.

Three design features may contribute to the shortfall of the learned rows. We did not test them separately, so the list names candidates, not measured causes.

The first candidate is the small size of most edit effects. Random-direction edits leave HellaSwag unchanged in every row, and their MMLU-Pro SD is 0.007. Abliteration's SDs are 0.013 and 0.005. Each target moves in steps of one item out of 100, so many edits differ by one item or not at all. Ranks over values like these are hard to estimate reliably.

The second candidate is item sampling in the targets. Each target is a 100-item accuracy, with a binomial standard error of 0.045 for MMLU-Pro and 0.049 for HellaSwag at the base fast-tier values (Section~\ref{sec:reliability}). Within family and device, the HellaSwag residual SD across the 300 D3 entries is 0.0131, no larger than the 0.0136 item standard error (\path{docs/council/2026-11-05/metrology_budget_output.txt}, lines 43 and 44). Item sampling therefore accounts for the residual spread that we can measure there. We have not estimated a reliability coefficient for either target, so we cannot say how much of the remaining rank disagreement comes from target noise.

The repeat data are narrower than a noise estimate. The base model and the ten control entries, re-evaluated on the same device with a fixed seed, matched exactly on every metric. Across hosts, the mean absolute difference on the 11 edited entries evaluated on the H100 was 0.008 on MMLU-Pro and 0.003 on HellaSwag, with Spearman 0.75 and 0.76 across those entries. On the A10, four edited entries gave mean differences of 0.010 and 0.003. Seed-to-seed SD on the fast tier is 0.0000 for both targets, because both are likelihood-scored (\path{docs/council/2026-11-05/metrology_budget_output.txt}, lines 24 and 25). These repeats bound host disagreement on a few entries. They do not measure the reliability of either target across the suite.

The third candidate is family and norm structure. The v2 design does not separate family from norm. Family means differ, and the delta-norm baseline reaches 0.21 and 0.18 on the trajectory split. Under family holdout its Spearman falls to $-$0.13 and $-$0.29. Section~\ref{sec:exploits} already noted that a pooled rank can reflect family offsets. Within-family refits on HellaSwag range from $-$0.12 to 0.71, so the signal is uneven across families. The suite does not adjust for family-by-norm structure.

Two further limits apply. The 260 targets with inferred runtime metadata block qualification (Section~\ref{sec:reliability}), and we did not test whether they change the result. The baselines are also single seed-0 fits, and their paired intervals condition on those fits.

The v2 evidence supports three narrower claims. On seed 0 of the trajectory split, a model with no weight statistics ranks edits nearly as well as StatNN: the Spearman gap is 0.042 on MMLU-Pro and 0.035 on HellaSwag. On the same split, StatNN has lower MAE than every learned row on both targets, and each paired interval excludes zero (Section~\ref{sec:v2}). The accuracy data also show an exploratory within-family signal in a residual analysis, which is not yet a confirmation.

With family, device and log-norm terms removed using training folds only, StatNN reduces residual absolute error relative to an isotonic-norm comparator by $\theta = 0.146$ [0.069, 0.214], over 345 rows and 278 clusters (\path{docs/council/2026-11-05/round2_residual_estimand_output.txt}, lines 2 to 11). Here $\theta = 1 - A/B$, where $A$ and $B$ are the summed absolute residuals of StatNN and of the comparator. The residual is built from MMLU-Pro accuracy on reused P4 data, and its specification was not frozen before we inspected the outcomes. The estimate is exploratory, and it does not confirm H1 under a new protocol.

The next study is a readability study on normalized gold-letter log-probability. It asks whether StatNN's residual advantage holds on fresh items and components, under rules fixed before any new outcome is scored. The primary outcome is the mean, over 1,000 paired items per edit, of the normalized log-probability of the gold answer letter, measured against the actual predecessor. The estimand is $\theta$, the row-weighted relative reduction in residual absolute error for fixed StatNN against a training-fitted isotonic-norm comparator, after training-only family, device and log-norm terms are removed. A pass needs $\theta \ge 0.15$, a lower 95\% bound above zero, a half-width no greater than $\theta/2$, at least 100 valid fresh components, measured execution qualification, and at least 80\% power at $\theta = 0.15$. The non-null stratum is fixed before outcomes as the family set {LoRA SFT, task vector, chain}. Zero-effect and destructive edits stay in the denominator. An upper 95\% bound below 0.05 on a qualified, adequately powered panel would kill the practical-readability claim for this endpoint and design only. Any other intermediate or wide result is inconclusive. The protocol is recorded in \path{docs/PREREG_H1_TF.md} (status: frozen prospective design, pending Board A acceptance and measured execution gates). If the gates fail, the contribution is a measurement and benchmark report with its negative results.
